\section{From failure documents to bowties: useful extraction without causal overclaiming}
\label{exec21:ace-scope}
A bank's operational-risk team may want to turn incident reports and control inventories into a diagram of causes, a critical event, preventive barriers, consequences and mitigating barriers. Such a representation can organize evidence and expose missing controls. It does not yet specify the effect of installing a control. ACE, by Venkatesh, Oudmans and Zurfluh, supplies a useful industrial analogue: three pipelines generate bowtie diagrams from Failure Mode and Effects Analysis documents. The 3 July 2026 publisher edition compares retrieval, OCR and a vision-assisted pipeline, with substantially better results on selected structured tables than on narrative text. Its central contribution for this proposal is a concrete extraction workflow and a set of failure modes to test before treating a diagram as a causal model.\citep{exec21_ace}

\subsection{What is extracted, and from which evidence?}
\label{exec21:ace-pipeline}
In the retrieval pipeline, Markdown FMEA content is chunked, embedded with all-MiniLM-L6-v2 and indexed with FAISS; part-name queries select context for a structured prompt. The OCR pipeline uses PaddleOCR and Img2Table to convert table images to Markdown before generation. In the vision pipeline, Mistral-Small produces a Markdown table and description for a second language model to convert into bowtie JSON. The reported implementation uses Python 3.9, sentence-transformers and local GGUF inference through llama-cpp-python. These are named dependencies and reported components, not a pinned, independently reproduced installation in this project.\citep{exec21_ace}

The evaluation first screens eight model variants on four causal-loop extraction texts and ten false-premise questions. It then selects three instruction-tuned model families and builds tabular references for a sensor, valve assembly and shaft from the US Navy reliability handbook. A domain expert constructs those reference diagrams. Narrative evaluation uses dynamic-seal text and three reference bowties. The paper explicitly notes that expert diagrams sometimes contain implicit variables absent from the source. This is a serious labeling distinction: agreement with an expert's augmented interpretation and faithful extraction of explicit text are different tasks.\citep{exec21_ace}

A bank-facing interface should therefore retain a source span or table-cell identity for each proposed role and edge. It should separately label an explicit statement, an expert-supplied hypothesis, a model inference and an unresolved conflict. If a narrative does not state a mitigating barrier or consequence, the result can be incomplete. A schema requiring every field to contain a plausible phrase would reward invention. ACE identifies precisely this problem in fallback prompts that ask for a mechanism and otherwise insert a generic mechanism label. Its narrative example also lacks enough information to complete the consequence side of the expert bowtie.\citep{exec21_ace}

The article's attribution license is recorded with the genuine PDF. It does not establish reuse rights for handbook data, model weights, pipeline code or dependencies. Recovering the actual prompts, preprocessing, model hashes, generation settings and evaluation files is a separate step before adopting an implementation. The proposal can adopt the representation and provenance requirements now while keeping executable reuse a conditional choice.

\subsection{Read metrics with their alignment and denominator}
\label{exec21:ace-metrics}
ACE matches predicted and reference labels with SBERT cosine similarity at an empirically selected threshold of 0.75. It reports node and edge precision, recall and F1 by role or edge type. Its reported graph edit distance is a count of false-positive and false-negative nodes and edges after matching, rather than the minimum edit path in the usual graph-theoretic problem. Threshold choice, matching constraints and reference roles are consequently part of the measurement contract. Semantically similar labels cannot be allowed to erase a reversal of causal direction or the distinction between a cause and a barrier.\citep{exec21_ace}

Table 4's selected configurations range from node F1 0.81 to 0.94 and edge F1 0.58 to 0.93. The paper describes a missed critical event yielding node F1 0.80 but edge F1 0.44 for a shaft case. On narratives, Table 5's average accuracy ratings range from 1.20 to 3.00 on a five-point scale. The strongest clarity rating, 3.80, is not evidence that the same output has correct causal structure. Expert-level records, rating units, missing ratings and aggregation rules remain necessary to interpret the reported means. These are small, selected industrial evaluations, not banking transfer results.\citep{exec21_ace}

The counts also deserve direct reconciliation. For each of two prompting conditions, hallucinated, unclear and correct-rejection responses are respectively $(21,10,49)$ and $(9,11,60)$, each summing to 80. Hallucinations decline from $21/80=26.25\%$ to $9/80=11.25\%$, a 15 percentage-point reduction. Unclear responses increase by one, despite prose describing them as reduced. Separately, 12 of 32 schema outputs are invalid JSON, or 37.5\%. These experiments have different denominators and cannot be merged into one accuracy measure. The reference calculation checks the printed arithmetic; no model responses are reproduced here.\citep{exec21_ace,exec21_ace_reference}

Aggregate F1 can conceal a consequential omission. Consider an authored chain of 101 nodes with 100 directed edges from node zero to node 100. A prediction retains every node and 99 correct edges but omits the bridge from node 49 to node 50. Node F1 is one, and
\begin{equation}
 F_{1,\mathrm{edge}}=\frac{2\,TP}{2\,TP+FP+FN}
 =\frac{198}{199}\simeq0.995.
 \label{exec21:ace-path}
\end{equation}
Nevertheless, the predicted graph contains no directed path from the initiating event to the final consequence. For a bank tracing a payment outage through failed controls to a customer loss, a critical-path or control-role error can outweigh many correctly extracted labels. This example motivates reporting decision-critical role and path errors alongside aggregate graph scores; it does not characterize an observed ACE run.\citep{exec21_ace_reference}

\subsection{Sensitivity indices do not establish interchangeability}
\label{exec21:ace-sensitivity}
The reported Sobol experiments vary prompt strictness, prompt type, context pipeline and selected decoding settings. They attribute much of the output variation to prompt factors and their interactions. The text interprets a near-zero first-order context index as evidence that the three pipelines are interchangeable under equivalent prompts. That conclusion needs more than a first-order variance decomposition.\citep{exec21_ace}

For independent design factors $P$ and $C$, let $f(P,C)$ be an evaluation score with positive variance. A first-order context index and a total-order context index are
\begin{equation}
 S_C=\frac{\operatorname{Var}(E[f\mid C])}{\operatorname{Var}(f)},
 \qquad
 S_{T,C}=\frac{E[\operatorname{Var}(f\mid P)]}{\operatorname{Var}(f)}.
 \label{exec21:ace-indices}
\end{equation}
These quantities concern variation under a specified experimental input distribution. To see the distinction directly, stipulate independent uniform binary $P,C$ and $f(P,C)=\mathbf 1\{P=C\}$. Both conditional means given $C$ are $1/2$, while the marginal variance is $1/4$. Thus $S_C=0$ but $S_{T,C}=1$. Changing context at fixed prompt changes the score by either one or minus one: the averaged main effect cancels, although the context matters completely through interaction. Exact enumeration verifies all four cells.\citep{exec21_ace_reference}

This witness does not show that ACE's actual total-order estimate is one. It establishes why a near-zero first-order estimate alone cannot prove equivalence. Even small estimated total effects over three component tables cannot establish agreement on new formats or rare critical failures. The next experiment should freeze categorical encodings, factor probabilities, semantic matching and decoding controls; retain the full sample matrix and all invalid outputs; and compare pipelines on paired held-out documents with prespecified tolerances for critical-role errors, ordinary extraction accuracy, latency and cost. The precise factor list in the expanded Sobol run, including whether decoding factors remain, should be recovered from code rather than inferred from its abbreviated paragraph.

ACE reports identical JSON in ten runs with different seeds for selected fixed configurations and attributes this to quantization and constrained prompts. Repeated identity establishes observed stability in those trials, not truth or its mechanism. A controlled precision comparison, verified seed propagation and decoding configuration, cache checks and raw token/output records would distinguish competing explanations. Ten repetitions of an erroneous diagram remain the same error. Schema-constrained prompting also should not be described as a hard grammar-constrained decoder without evidence of such an implementation.\citep{exec21_ace}

\subsection{From a role diagram to an intervention question}
\label{exec21:ace-intervention}
An unsigned arrow is insufficient to price an intervention. In a deliberately minimal example, both structural laws $Y=A$ and $Y=1-A$ have topology $A\to Y$, with binary action $A$ and loss event $Y$. Their intervention contrasts $E[Y\mid do(A=1)]-E[Y\mid do(A=0)]$ are respectively one and minus one. If a loss event costs an assumed USD 10 and the action costs USD 2, the incremental margin is
\begin{equation}
 V=-10\{E[Y\mid do(A=1)]-E[Y\mid do(A=0)]\}-2,
 \label{exec21:ace-value}
\end{equation}
giving USD $-12$ or USD 8. These two laws illustrate missing mechanism information in a common topology; they are not claimed to have identical observational distributions or to satisfy the same signed barrier interpretation. A diagram extractor has not supplied either law merely by naming an edge.\citep{exec21_ace_reference}

For a bank, the translation layer must additionally identify the action that can be changed, the time at which it is assigned, the eligible unit and population, the loss horizon, the data-generating and observation processes, and any confounding or feedback. A generic barrier node might represent an existing control, a proposed policy, a measured proxy or an expert belief. These possibilities imply different models. Time indexing may be necessary before a feedback description can be represented as an acyclic causal query. Unresolved semantics should remain explicit and should prevent an automatic effect estimate.

\subsection{A bounded agentic workflow and public validation route}
\label{exec21:ace-integration}
ACE proposes future agentic correction and knowledge-graph reasoning, while its reported pipelines are single-pass and single-component. The proposal should keep that boundary visible. A candidate workflow retrieves immutable evidence, proposes typed roles and edges, checks schema and source support, and asks a bounded correction stage to address a specific omission or contradiction. It preserves all attempts and unresolved results. A graph database can store those objects, but neither storage nor repeated agent agreement validates a causal claim.\citep{exec21_ace}

The CLI or MCP contract should expose source hashes and spans, role types, temporal scope, explicit versus inferred status, schema/model/prompt versions, costs and distinct failure states. Invalid JSON, missing source evidence, a disputed expert addition and a scientifically unidentified intervention are different outcomes. No pipeline should convert all four into a generic failure, fabricate a completed bowtie, or silently grant an intervention-analysis capability.

Public documents and generated mechanisms can justify the next experiment before costly institutional evidence is requested. Split by document family or component, not repeated prompts on the same case. Include explicit and implicit causes, absent barriers, contradictory statements, multi-line tables and misleading nearby components. Freeze the reference construction and separate text-supported truth from expert augmentation. Report invalid outputs, critical-role/path loss, false additions, disagreement, correction attempts and total review cost. Add known-mechanism tasks only where the desired intervention is defined, and retain refusal when extraction cannot support it. Such a prototype can establish whether assisted diagram construction merits a bank trial without claiming that the trial's outcome is already known.\citep{exec21_ace,exec21_ace_reference}
